\documentclass[11pt]{scrartcl}

\usepackage[stixtwo]{fontsetup}
\usepackage{amsmath}

\begin{document}
Mivel $x^2 - 2x + 1 = (x - 1)^2$ és $x^3 - x = x(x - 1)(x + 1)$, ezért
\[
	f(x) = \frac{|x - 1| \sqrt{x^4 + 3}}{x (x - 1)(x + 1)}.
\]
\begin{itemize}
	\item $x < 1$ esetén $|x - 1| = 1 - x$, így
	      \[
		      \lim_{x \to -\infty} f(x)
		      = \lim_{x \to -\infty}
		      \frac{1 - x}{x} \cdot
		      \frac{\sqrt{1 + \frac{3}{x^4}}}{1 - \frac{1}{x^2}}
		      = -1 \cdot 1
		      = -1.
	      \]

	\item \[
		      \lim_{x \to 1 \pm 0} f(x)
		      = \lim_{x \to 1 \pm 0}
		      \underbrace{\frac{|x - 1|}{x - 1}}_{= \pm 1} \cdot
		      \underbrace{\frac{\sqrt{x^4 + 3}}{x (x + 1)}}_{\to 1}
		      = \pm 1,
	      \]
	      a bal és a jobb oldali határérték különbözik, így
	      $\lim_{x \to 1} f(x)$ nem létezik.

	\item $f$ folytonos a $2$ pontban, így
	      \(
	      \lim_{x \to 2} f(x) = f(2) = \frac{\sqrt{19}}{6}
	      \).
\end{itemize}
\end{document}
